\documentclass[12pt,letterpaper]{article}



%%
%% Required packages:
%%
\usepackage{etex}
\usepackage{amsmath}
\usepackage{amssymb}
\usepackage{amstext}
\usepackage{amsthm}
\usepackage{fancyhdr,fancybox}
\usepackage{mathrsfs} % need for \mathscr command
\usepackage{multicol}
\usepackage[normalem]{ulem}
%\usepackage{pictex,pictexplus}
% \usepackage{pictexwd}
\usepackage{rotating}
\usepackage{url}
\usepackage{xfrac}
\usepackage{booktabs}
\usepackage{color,soul}
\usepackage{comment}

\usepackage{hyperref}
\hypersetup{
    colorlinks=true,     % false: boxed links; true: colored links
    linkcolor=blue,      % color of internal links
    citecolor=blue,      % color of links to bibliography
    filecolor=blue,      % color of file links
    urlcolor=blue        % color of external links
}


\usepackage{tikz}
\usetikzlibrary{backgrounds,calc}

%%
%% Common formatting commands
%%
\setlength{\parindent}{0in}
\setlength{\textwidth}{7in}
\setlength{\evensidemargin}{-0.25in}
\setlength{\oddsidemargin}{-0.25in}
\setlength{\parskip}{.5\baselineskip}
\setlength{\topmargin}{-0.5in}
\setlength{\textheight}{9in}

%               Problem and Part
\newcounter{problemnumber}
\newcounter{partnumber}
\newcounter{subpartnumber}
\newcommand{\Problem}{%
  \refstepcounter{problemnumber}%
  \setcounter{partnumber}{0}%
  \setcounter{subpartnumber}{0}%
  \item[\makebox{\hfil\textbf{\theproblemnumber.}\hfil}]%
  }
\newcommand{\InProblem}[2][1.6]{%
  \refstepcounter{problemnumber}%
  \setcounter{partnumber}{0}%
  \setcounter{subpartnumber}{0}%
  \textbf{\theproblemnumber.}\ \ \parbox[t]{#1in}{#2}%
}
\newcommand{\InSmallProblem}[1]{\InProblem[1.05]{#1}}
\newcommand{\NoProblem}[1]{%
  \mbox{\hfil\textbf{#1}\hfil}%
  }
\newcommand{\UnProblem}{%
  \refstepcounter{problemnumber}%
  \setcounter{partnumber}{0}%
  \setcounter{subpartnumber}{0}%
  \mbox{\hfil\textbf{\theproblemnumber.}\hfil}%
  }
\newcommand{\InFlexProblem}[1]{%
  \refstepcounter{problemnumber}%
  \setcounter{partnumber}{0}%
  \setcounter{subpartnumber}{0}%
  \textbf{\theproblemnumber.}\ \ #1%
}
%%  Part:
\newcommand{\Part}{%
  \renewcommand{\thepartnumber}{(\alph{partnumber})}%
  \refstepcounter{partnumber}
  \setcounter{subpartnumber}{0}%
  \item[(\alph{partnumber})]%
}
\newcommand{\InPart}[2][1.75]{%
  \renewcommand{\thepartnumber}{(\alph{partnumber})}%
  \refstepcounter{partnumber}%
  \setcounter{subpartnumber}{0}%
  (\alph{partnumber})\ \ \parbox[t]{#1in}{#2}%
}
\newcommand{\UnPart}{%
  \refstepcounter{partnumber}%
  \renewcommand{\thepartnumber}{(\alph{partnumber})}%
  \setcounter{subpartnumber}{0}%
  (\alph{partnumber})%
}
\newcommand{\InLargePart}[1]{\InPart[3.05]{#1}}
\newcommand{\InFullPart}[1]{\InPart[6.2]{#1}}
\newcommand{\InSmallPart}[1]{\InPart[1.05]{#1}}
%% SubPart:
\newcommand{\SubPart}{%
  \refstepcounter{subpartnumber}%
  \item[(\roman{subpartnumber})]%
  }
\newcommand{\InSubPart}[2][2.25]{%
  \stepcounter{subpartnumber}%
  (\roman{subpartnumber})\ \ \parbox[t]{#1in}{#2}%
}
\newcommand{\InSmallSubPart}[1]{\InSubPart[1.5]{#1}}
\newcommand{\InTinySubPart}[1]{%
  \stepcounter{subpartnumber}%
  (\roman{subpartnumber})\ \ \makebox{#1}%
}
\newcommand{\InMySubPart}[2][2.25]{%
  \stepcounter{subpartnumber}%
  \parbox[t]{#1in}{\makebox[0.3in][l]{(\roman{subpartnumber})}\ \ {#2}}%
}
\newcommand{\TrueFalse}{\textbf{T}\ \ \textbf{F}\ \ }
\newcommand{\TFPart}[1]{% 
  \stepcounter{partnumber}%
  \setcounter{subpartnumber}{0}%
  \item[(\alph{partnumber})] \TrueFalse\parbox[t]{5.5in}{#1}}

\newcommand{\Points}[1]{(#1 points) \ }
\newcommand{\Point}[1]{(#1 point) \ }


%% For Answers:
\newcounter{answernumber}
\newcommand{\Answer}{%
  \refstepcounter{answernumber}%
  \setcounter{partnumber}{0}%
  \setcounter{subpartnumber}{0}%
  \item[\makebox{\hfil\textbf{\theanswernumber.}\hfil}]%
  }


% Setting up the headers for the first page
%\chead{} % will default to what is typed in the file.
\fancyhead[L]{\textsc{Math 22a}}
\fancyhead[R]{\textsc{Fall 2024}}
\fancyfoot[R]{
	\vspace*{\fill}\footnotesize\textcopyright{} \emph{Harvard Math 22a}	
}
\renewcommand{\headrulewidth}{0pt}

%\fancyhead[C]{}
%	\chead{}	
%	\lhead{}
%	\rhead{}
%	\cfoot{}


% Putting copyright on the footer of each page
%\fancypagestyle{secondpage}{
%	\fancyhead[C]{TESTing again} % change this to be blank
%		%\chead{}	
%	\fancyhead[L]{bears} % change this to be blank
%	\fancyhead[R]{dogs} % change this to be blank
%	\fancyfoot[R]{ %keeps the footer the same
%		\vspace*{\fill}\footnotesize\textcopyright{} \emph{Harvard Math 22a}	
%	}
%}
%\renewcommand{\headrulewidth}{0pt}
%\pagestyle{fancy}
\pagestyle{empty}


%\lhead{}
%\rhead{}
%\rfoot{\vspace*{\fill}\footnotesize\textcopyright{} \emph{Harvard Math 22a}}
%\renewcommand{\headrulewidth}{0pt}
%\pagestyle{fancy}




%%
%% Common shortcut commands:
%%
\newcommand{\ds}{\displaystyle}
\newcommand{\sgn}{\operatorname{sgn}}
\renewcommand{\Pr}{\operatorname{P}}
\newcommand{\CPr}[3][{}]{\Pr_{\text{#1}}\left(#2\,|\,#3\right)}

\newcommand{\C}{\mathbb{C}}
\newcommand{\R}{\mathbb{R}}
\newcommand{\Z}{\mathbb{Z}}
\newcommand{\N}{\mathbb{N}}
\newcommand{\Q}{\mathbb{Q}}
\newcommand{\D}{\mathbb{D}}

\newcommand{\rref}{\operatorname{rref}}
\newcommand{\im}{\operatorname{im}}
\newcommand{\rank}{\operatorname{rank}}
\newcommand{\diag}{\operatorname{diag}}
\newcommand{\Span}{\operatorname{span}}
%\renewcommand{\vec}[1]{\mathbf{#1}}
\newcommand{\charpoly}{\mathscr{P}}
\newcommand{\nicefrac}[2]{\sfrac{#1}{#2}} % using xfrac package
\newcommand{\topology}{\mathcal{T}}
\newcommand{\Inn}{\operatorname{Inn}}

\newcommand{\blankbox}[2]{\fbox{\rule{#1}{0in}\rule{0in}{#2}}}

\newenvironment{myindentpar}[1]%
 {\begin{list}{}%
         {\setlength{\leftmargin}{#1}
          \setlength{\rightmargin}{\leftmargin}}%
         \item[]%
 }
 {\end{list}}
\newtheorem{theorem}{Theorem}
\newtheorem*{NStheorem}{Theorem}
\newtheorem{cor}[theorem]{Corollary}
\newtheorem*{claim}{Claim}
\newtheorem{defn}{Definition}

\newenvironment{amatrix}[1]{%
  \left(\begin{array}{@{}*{#1}{c}|c@{}}
}{%
  \end{array}\right)
}

%--------grstep
% For denoting a Gauss' reduction step.
% Use as: \grstep{\rho_1+\rho_3} or \grstep[2\rho_5 \\ 3\rho_6]{\rho_1+\rho_3}
\newcommand{\grstep}[2][\relax]{%
   \ensuremath{\mathrel{
       {\mathop{\longrightarrow}\limits^{#2\mathstrut}_{
                                     \begin{subarray}{l} #1 \end{subarray}}}}}}
\newcommand{\swap}{\leftrightarrow}



\usepackage{logicpuzzle}


\chead{\Large\textbf{Puzzles and Linear Systems, $ \vec \S$1.1}}% 

\begin{document}
\thispagestyle{fancy}

\begin{enumerate}
  \Problem A Kakuro(-ish) Puzzle. Enter numbers from 1 to 9 in any order into the blank cells. The rows and columns should sum to the number on the grid. Within a given row or column, no number can be repeated.
  
  \definecolor{kakuro}{RGB}{155,206,167}
  \kakurosetup{color=kakuro,rows=4,columns=4}
  \begin{center}
  	\begin{kakuro}
	\kakurorow{4}{\Black,\KKR{10}{},\KKR{13}{},\Black}
	\kakurorow{3}{\KKR{}{17},9,1,\KKR{13}{}}
	\kakurorow{2}{\KKR{}{7},6,5,3}
	\kakurorow{1}{\Black,\KKR{}{12},3,2}
%	\kakurorow{1}{\Black,\Black,\Black,\Black,\Black}
\end{kakuro}
%  	\hspace{1.5cm}
%  	\begin{kakuro}[solution]
%  		\framepuzzle
%  		\kakurorow{5}{\Black,\KKR{23}{},\KKR{16}{},\KKR{10}{},\Black}
%  		\kakurorow{4}{\KKR{}{14},9,1,4,\KKR{3}{}}
%  		\kakurorow{3}{\KKR{}{16},6,5,3,2}
%  		\kakurorow{2}{\KKR{}{14},8,3,2,1}
%  		\kakurorow{1}{\Black,\KKR{}{8},7,1,\Black}
%  	\end{kakuro}
  
  \end{center}
This puzzle is taken from ``Linear and Complex Analysis for Applications'' by John D.Angelo.


\newpage

\begin{defn}
A {\bf linear equation} in the variables $x_1, \dots, x_n$ is an equation that can be put in the following form
\begin{equation*}
a_1 x_1 +a_2 x_2 +\dots + a_n x_n =b
\end{equation*}
where $b$ and the coefficients $a_1,\dots, a_n$ are real or complex numbers. A {\bf system of linear equations} is a collection of one or more linear equations involving the same variables. A {\bf solution} to the linear system is a list of $(s_1, \dots, s_n)$ of numbers that are solutions to each of the linear equations in the system. 
\end{defn}

Given a linear system, some fundamental questions are:
\begin{itemize}
\item Does there exist at least one solution?
\item If a solution exists, how can we find it?
\item If a solution exists, is it unique?
\end{itemize}

\begin{defn}
A linear system is {\bf consistent} if there is a solution and {\bf inconsistent} if there is no solution.
\end{defn}

\Problem Find examples of the following.
\Part An inconsistent linear system in two variables with two equations.
%\begin{flalign}
%y-x& =0 && \nonumber \\
%y-x  &=1 && \nonumber  
%\end{flalign}

\vfill

\Part A consistent linear system in two variables with two equations.
%\begin{flalign}
%y-x& =0 && \nonumber \\
%2y-2x  &=0 && \nonumber  
%\end{flalign}

\vfill

%\Part
%\begin{flalign}
%y-x& =0 && \nonumber \\
%y+x  &=0 && \nonumber  
%\end{flalign}

%\vfill

\newpage

{\bf Fact:} A linear system has either no solution, exactly one solution, or infinitely many solutions.

\begin{defn}
The following elementary row operations preserve solutions.
\begin{enumerate}
\item Replace one row by the sum of itself and a multiple of another row.
\item Interchange two rows.
\item Multiply all entries in a row by a non-zero constant.
\end{enumerate}
\end{defn}

\begin{defn}
Two matrices are called {\bf row equivalent} if there is a sequence of elementary row operations changing one matrix to another.
\end{defn}

{\bf Fact:} If the augmented matrices of two linear systems are row equivalent, then they both have the same solution set. 

\Problem Solve the following systems.

\Part 
\begin{flalign}
2x+3y& =0 && \nonumber \\
4x+5y  &=1 && \nonumber  
\end{flalign}

\vfill

\Part 
\begin{flalign}
2x+3y& =0 && \nonumber \\
4x+6y  &=1 && \nonumber  
\end{flalign}

\vfill


\newpage

\Problem Lights Out Puzzle. The yellow squares represent rooms with the lights on, and the white squares represent squares with lights off. The goal is to turn off all the lights, but there is a catch. When you flip the switch in any room, the lights in the rooms horizontally and vertically adjacent are also flipped. 
%http://mathworld.wolfram.com/LightsOutPuzzle.html
\begin{center}
\begin{chaossudoku}[rows=3,columns=3]
   % \setrow{5}{4,3,5,1,2}
  %  \setrow{4}{2,1,3,4,5}
 %   \setrow{3}{,,}
  %  \setrow{2}{,,}
  %  \setrow{1}{,,}
    \begin{puzzlebackground}
%\fillarea{Wheat}{(1,1)--(1,2)--(2,2)--(2,3)--(4,3)--(4,1) --(1,1)}
\fillarea{yellow}{(2,3)--(2,4)--(3,4)--(3,3)--(2,3)}
\fillarea{yellow}{(2,2)--(2,3)--(3,3)--(3,2)--(2,2)}
\fillarea{yellow}{(2,1)--(2,2)--(3,2)--(3,1)--(2,1)}

\fillarea{yellow}{(1,2)--(1,3)--(2,3)--(2,2)--(1,2)}

\fillarea{yellow}{(3,1)--(3,2)--(4,2)--(4,1)--(3,1)}
%\fillarea{yellow}{(2,3)--(2,5)--(3,5)--(3,4)--(5,4) --(5,2)--(4,2)--(4,3)--(2,3)}
%\fillarea{blue}{(3,4)--(3,6)--(6,6)--(6,5)--(5,5) --(5,4)--(3,4)}
%\fillarea{yellow}{(4,1)--(4,2)--(5,2)--(5,5)--(6,5) --(6,1)--(4,1)}
    \end{puzzlebackground}
  \end{chaossudoku}
\end{center}

\vfill

\Problem Write a short summary of today's key ideas.

\vfill

\end{enumerate}

%\end{document}

\newpage
\chead{\Large\textbf{Puzzles and Systems, $\vec \S$1.1 -- Answers / Solutions}}%
\lhead{}
\rhead{}
\thispagestyle{fancy}

\begin{enumerate}
  \Answer Here are the two possible solutions:
  
   \definecolor{kakuro}{RGB}{155,206,167}
  \kakurosetup{color=kakuro,rows=4,columns=4}
  \begin{center}
  	\begin{kakuro}[solution]
	\kakurorow{4}{\Black,\KKR{10}{},\KKR{13}{},\Black}
	\kakurorow{3}{\KKR{}{17},9,8,\KKR{13}{}}
	\kakurorow{2}{\KKR{}{7},1,2,4}
	\kakurorow{1}{\Black,\KKR{}{12},3,9}
%	\kakurorow{1}{\Black,\Black,\Black,\Black,\Black}
\end{kakuro}
  	\hspace{1.5cm}
  	\begin{kakuro}[solution]
	\kakurorow{4}{\Black,\KKR{10}{},\KKR{13}{},\Black}
	\kakurorow{3}{\KKR{}{17},8,9,\KKR{13}{}}
	\kakurorow{2}{\KKR{}{7},2,1,4}
	\kakurorow{1}{\Black,\KKR{}{12},3,9}
%	\kakurorow{1}{\Black,\Black,\Black,\Black,\Black}
\end{kakuro}
  \end{center}
  
    This puzzle is taken from ``Linear and Complex Analysis for Applications" by John D'Angelo. 
  
  \Answer
  \Part An inconsistent linear system in two variables with two equations.
\begin{flalign}
y-x& =0 && \nonumber \\
y-x  &=1 && \nonumber  
\end{flalign}

\Part A consistent linear system in two variables with two equations with infinitely many solutions.
\begin{flalign}
y-x& =0 && \nonumber \\
2y-2x  &=0 && \nonumber  
\end{flalign}

\Part A consistent linear system in two variables with two equations with exactly one solution.
\begin{flalign}
y-x& =0 && \nonumber \\
y+x  &=0 && \nonumber  
\end{flalign}

%\vfill

\Answer 

\Part We begin by translating the linear system into the corresponding augmented matrix:

  $$  \begin{amatrix}{2}
    2&3&0\\%
    4&5&1
    \end{amatrix}$$
 
Now we can perform elementary row operations to simplify the matrix:

  
    \[
\begin{amatrix}{2}
    2&3&0\\%
    4&5&1
    \end{amatrix}
    %
    \grstep{R_2 - 2 R_1}
    %
\begin{amatrix}{2}
    2&3&0\\%
    0&-1&1
    \end{amatrix} 
        \grstep{ -  R_2}
\begin{amatrix}{2}
    2&3&0\\%
    0&1&-1
    \end{amatrix}
        \grstep{R_1 - 3 R_2}
\begin{amatrix}{2}
    2&0&3\\%
    0&1&-1
    \end{amatrix}
            \grstep{\frac{1}{2}R_1}
\begin{amatrix}{2}
    1&0&\frac{3}{2}\\%
    0&1&-1
    \end{amatrix},
    \]
 which becomes the following system of equations:
 \begin{align*}
x& =\frac{3}{2}  \\
y  &=-1   
\end{align*}
\Part We begin again by translating the linear system into the corresponding augmented matrix:

  $$  \begin{amatrix}{2}
    2&3&0\\%
    4&6&1
    \end{amatrix}$$
 
Now we can perform elementary row operations to simplify the matrix:

  
    \[
\begin{amatrix}{2}
    2&3&0\\%
    4&6&1
    \end{amatrix}
    %
    \grstep{R_2 - 2 R_1}
    %
\begin{amatrix}{2}
    2&3&0\\%
    0&0&1
    \end{amatrix} ,
    \]
 which becomes the following system of equations:
 \begin{align*}
2x+3y& =0  \\
0 &=1.   
\end{align*}

Thus the system has no solutions and it is inconsistent. Geometrically, we can see that the equations are parallel lines.

  
  \Answer We start in the upper-left corner:   
   \begin{center}
\begin{chaossudoku}[rows=3,columns=3]
   % \setrow{5}{4,3,5,1,2}
  %  \setrow{4}{2,1,3,4,5}
 %   \setrow{3}{,,}
  %  \setrow{2}{,,}
  %  \setrow{1}{,,}
    \begin{puzzlebackground}
%\fillarea{Wheat}{(1,1)--(1,2)--(2,2)--(2,3)--(4,3)--(4,1) --(1,1)}
%\fillarea{yellow}{(2,3)--(2,4)--(3,4)--(3,3)--(2,3)}
\fillarea{yellow}{(2,2)--(2,3)--(3,3)--(3,2)--(2,2)}
\fillarea{yellow}{(2,1)--(2,2)--(3,2)--(3,1)--(2,1)}

%\fillarea{yellow}{(1,2)--(1,3)--(2,3)--(2,2)--(1,2)}

\fillarea{yellow}{(1,3)--(1,4)--(2,4)--(2,3)--(1,3)}

\fillarea{yellow}{(3,1)--(3,2)--(4,2)--(4,1)--(3,1)}
%\fillarea{yellow}{(2,3)--(2,5)--(3,5)--(3,4)--(5,4) --(5,2)--(4,2)--(4,3)--(2,3)}
%\fillarea{blue}{(3,4)--(3,6)--(6,6)--(6,5)--(5,5) --(5,4)--(3,4)}
%\fillarea{yellow}{(4,1)--(4,2)--(5,2)--(5,5)--(6,5) --(6,1)--(4,1)}
    \end{puzzlebackground}
  \end{chaossudoku}
\end{center}
Next we turn the upper middle light to get the following:
   \begin{center}
\begin{chaossudoku}[rows=3,columns=3]
   % \setrow{5}{4,3,5,1,2}
  %  \setrow{4}{2,1,3,4,5}
 %   \setrow{3}{,,}
  %  \setrow{2}{,,}
  %  \setrow{1}{,,}
    \begin{puzzlebackground}
%\fillarea{Wheat}{(1,1)--(1,2)--(2,2)--(2,3)--(4,3)--(4,1) --(1,1)}
%\fillarea{yellow}{(2,3)--(2,4)--(3,4)--(3,3)--(2,3)}
%\fillarea{yellow}{(2,2)--(2,3)--(3,3)--(3,2)--(2,2)}
\fillarea{yellow}{(2,1)--(2,2)--(3,2)--(3,1)--(2,1)}

%\fillarea{yellow}{(1,2)--(1,3)--(2,3)--(2,2)--(1,2)}
\fillarea{yellow}{(2,3)--(2,4)--(3,4)--(3,3)--(2,3)}

\fillarea{yellow}{(3,3)--(3,4)--(4,4)--(4,3)--(3,3)}

%\fillarea{yellow}{(1,3)--(1,4)--(2,4)--(2,3)--(1,3)}

\fillarea{yellow}{(3,1)--(3,2)--(4,2)--(4,1)--(3,1)}
%\fillarea{yellow}{(2,3)--(2,5)--(3,5)--(3,4)--(5,4) --(5,2)--(4,2)--(4,3)--(2,3)}
%\fillarea{blue}{(3,4)--(3,6)--(6,6)--(6,5)--(5,5) --(5,4)--(3,4)}
%\fillarea{yellow}{(4,1)--(4,2)--(5,2)--(5,5)--(6,5) --(6,1)--(4,1)}
    \end{puzzlebackground}
  \end{chaossudoku}
\end{center}

Now we turn the light in the upper right to get the following:

   \begin{center}
\begin{chaossudoku}[rows=3,columns=3]
   % \setrow{5}{4,3,5,1,2}
  %  \setrow{4}{2,1,3,4,5}
 %   \setrow{3}{,,}
  %  \setrow{2}{,,}
  %  \setrow{1}{,,}
    \begin{puzzlebackground}
%\fillarea{Wheat}{(1,1)--(1,2)--(2,2)--(2,3)--(4,3)--(4,1) --(1,1)}
%\fillarea{yellow}{(2,3)--(2,4)--(3,4)--(3,3)--(2,3)}
%\fillarea{yellow}{(2,2)--(2,3)--(3,3)--(3,2)--(2,2)}
\fillarea{yellow}{(2,1)--(2,2)--(3,2)--(3,1)--(2,1)}

%\fillarea{yellow}{(1,2)--(1,3)--(2,3)--(2,2)--(1,2)}
%\fillarea{yellow}{(2,3)--(2,4)--(3,4)--(3,3)--(2,3)}

\fillarea{yellow}{(3,2)--(3,3)--(4,3)--(4,2)--(3,2)}

%\fillarea{yellow}{(1,3)--(1,4)--(2,4)--(2,3)--(1,3)}

\fillarea{yellow}{(3,1)--(3,2)--(4,2)--(4,1)--(3,1)}
%\fillarea{yellow}{(2,3)--(2,5)--(3,5)--(3,4)--(5,4) --(5,2)--(4,2)--(4,3)--(2,3)}
%\fillarea{blue}{(3,4)--(3,6)--(6,6)--(6,5)--(5,5) --(5,4)--(3,4)}
%\fillarea{yellow}{(4,1)--(4,2)--(5,2)--(5,5)--(6,5) --(6,1)--(4,1)}
    \end{puzzlebackground}
  \end{chaossudoku}
\end{center}

Now going in the lower right turns off all the lights.

\end{enumerate}



\end{document}


%%% Local Variables: 
%%% mode: latex
%%% TeX-master: t
%%% End: 
